\documentclass[11pt,a4paper]{article}
\usepackage{amsmath,amssymb}
\newcommand{\Tr}{\operatorname{Tr}}
\begin{document}
\title{ALF entropy of a tracial binary shift}
\author{}
\date{}
\maketitle
\section*{Setup and source}
Consider the bilateral binary-shift algebra generated by self-adjoint unitaries $u_j$, $j\in\mathbb Z$, with
\[u_iu_j=(-1)^{b_{|i-j|}}u_ju_i,\qquad\alpha(u_j)=u_{j+1}.\]
Here $b_k\in\{0,1\}$ for $k\geq1$. Use the canonical invariant trace $\tau$, which vanishes on every nonempty ordered monomial with distinct indices. These are the conventions in E.~St\o rmer, \emph{A survey of noncommutative dynamical entropy}, arXiv:math/0007010 (2000), Section 9, pp.~31--33. Theorem 9.3 establishes existence of binary shifts with zero CNT entropy. The specific page 219 of Neshveyev and St\o rmer's book was not checked here.

\section*{Exact history calculation}
An operational partition $X=(x_a)$ satisfies $\sum_ax_a^*x_a=I$. Fix the history convention
\[K_{\boldsymbol a}=\alpha^{n-1}(x_{a_{n-1}})\cdots\alpha(x_{a_1})x_{a_0}.\]
Its ALF correlation matrix and entropy are
\[\rho^{(n)}_{\boldsymbol a,\boldsymbol c}=\tau(K_{\boldsymbol c}^*K_{\boldsymbol a}),\qquad S_n=-\Tr(\rho^{(n)}\log\rho^{(n)}).\]
Transposing the matrix to use the opposite Gram-matrix indexing convention leaves its spectrum unchanged. All logarithms are natural. Choose
\[x_0=I/\sqrt2,\qquad x_1=u_0/\sqrt2.\]
Then $x_0^*x_0+x_1^*x_1=I$, and
\[K_{\boldsymbol a}=2^{-n/2}u_{n-1}^{a_{n-1}}\cdots u_0^{a_0},\qquad\boldsymbol a\in\{0,1\}^n.\]
For equal strings the product $K_{\boldsymbol a}^*K_{\boldsymbol a}$ is $2^{-n}I$. For unequal strings, commutation and cancellation leave, up to a sign, a nonempty monomial supported on the symmetric difference of the strings. Its trace vanishes. Consequently,
\[\rho^{(n)}=2^{-n}I_{2^n},\qquad S_n=n\log2,\qquad h_\tau(\alpha,X)=\log2.\]
This proves a positive ALF lower bound for every admissible partition class containing $X$, regardless of the bitstream. It is an exact infinite-system statement, not a finite-size extrapolation. The projective partition $((I+u_0)/2,(I-u_0)/2)$ is obtained from $X$ by a two-dimensional Hadamard unitary on the outcome index. Its refined histories are therefore related by the tensor power of that unitary and have the same correlation-matrix spectrum.

\section*{Exact value for strictly local partitions}
Let $\mathcal A_{\rm loc}$ be the union of the finite interval algebras generated by the $u_j$. Define $h_{\rm ALF}^{\rm loc}$ as the supremum of $\limsup_n S_n/n$ over finite operational partitions in this algebra. Any such partition is supported in a fixed interval $[r,s]$ of length $L=s-r+1$. Every length-$n$ history belongs to the interval algebra on $[r,s+n-1]$. Ordered monomials span this algebra, whose vector-space dimension is at most $2^{n+L-1}$. Since the history matrix is a Gram matrix in $L^2(\tau)$,
\[\operatorname{rank}\rho^{(n)}\leq2^{n+L-1},\qquad S_n\leq(n+L-1)\log2.\]
Dividing by $n$, taking the upper limit, and using the explicit partition above gives
\[h_{\rm ALF}^{\rm loc}(\alpha,\tau)=\log2.\]
No approximation theorem extending this upper bound to arbitrary nonlocal partitions is asserted here. For the full von Neumann algebra the calculation proves at least $h_{\rm ALF}\geq\log2$.

\section*{Interpretation and limitations}
In particular the standard tracial binary shifts with zero CNT entropy are not counterexamples to positivity of ALF entropy. This does not prove that the algebraic K-property implies ALF positivity for general systems: the argument used the special orthogonality of binary monomials. Moreover, even an asymptotic inequality $h_{\rm ALF}\geq h_{\rm CNT}$ would give only the trivial lower bound zero when CNT vanishes. A one-step entropy comparison alone does not establish an asymptotic rate comparison. No inference concerning quantum Lyapunov exponents is made.
\end{document}
